%-------------------------------------------------------------------------------
%	SECTION TITLE
%-------------------------------------------------------------------------------
\cvsection{منتخب پروژه‌ها}


%-------------------------------------------------------------------------------
%	CONTENT
%-------------------------------------------------------------------------------

\begin{cventries}
	
	\cventry
	{پلتفرم \lr{Agentic} چندمرحله‌ای برای تولید داستان بر اساس اساطیر روم}
	{\lr{Python, LangChain, LLM Agents}}
	{پروژه شخصی}
	{سال 1403} % پیشنهاد: درج سال
	{
		\begin{cvitems}
			\item {طراحی و پیاده‌سازی یک جریان چندمرحله‌ای مبتنی بر \lr{LLM Agent} برای ایده‌پردازی، شخصیت‌پردازی، طراحی پیرنگ و تولید متن داستانی.}
			\item {به‌کارگیری \lr{LangChain} برای مدیریت زنجیره پردازش، طراحی \lr{Prompt}های مرحله‌ای و کنترل ساختار خروجی.}
			\item {بهینه‌سازی فرایند تولید محتوای روایی با رویکرد ماژولار جهت افزایش انسجام داستان و قابلیت توسعه‌پذیری.}
		\end{cvitems}
	}
	
	\cventry
	{پلتفرم استخراج ساختارمند داده از مقالات علمی حوزه نانوذرات پزشکی}
	{\lr{Python, LangExtract, JSON Schema, Scientific NLP}}
	{پروژه پژوهشی}
	{سال 1404}
	{
		\begin{cvitems}
			\item {طراحی \lr{pipeline} استخراج اطلاعات از مقالات علمی و آزمایشگاهی در حوزه نانوذرات پزشکی و کاربردهای درمان سرطان.}
			\item {پیش‌پردازش و پاک‌سازی متن مقالات، مهندسی ویژگی‌ها و استخراج داده‌های کلیدی در قالب ساختارمند \lr{JSON}.}
			\item {توسعه ساختار خروجی استاندارد با استفاده از \lr{LangExtract} جهت آماده‌سازی داده‌ها برای تحلیل و مستندسازی.}
		\end{cvitems}
	}
	
	\cventry
	{پلتفرم \lr{OCR} برای تبدیل مقالات علمی \lr{PDF} به متن و \lr{Markdown}}
	{\lr{Python, FastAPI, OCR, Document Processing}}
	{پروژه پژوهشی}
	{سال 1404}
	{
		\begin{cvitems}
			\item {پیاده‌سازی سرویس پردازش اسناد علمی مبتنی بر مدل \lr{OlmOCR} جهت استخراج دقیق متن از فایل‌های \lr{PDF}.}
			\item {حذف \lr{metadata} و بخش‌های زائد مقاله، بازسازی متن و تبدیل محتوای علمی به فرمت‌های متنی و \lr{Markdown}.}
			\item {طراحی فرایند آماده‌سازی اسناد علمی جهت استفاده در تحلیل متن، بازیابی اطلاعات و پردازش‌های مبتنی بر \lr{LLM}.}
		\end{cvitems}
	}
	
	\cventry
	{سرور اجرای مدل تولید تصویر \lr{FLUX.1-dev}}
	{\lr{Python, FastAPI, Modal, Lightning AI, FLUX.1-dev}}
	{پروژه شخصی}
	{سال 1403}
	{
		\begin{cvitems}
			\item {توسعه \lr{API Server} با \lr{FastAPI} برای اجرای مدل مولد تصویر \lr{FLUX.1-dev} روی زیرساخت‌های \lr{Modal} و \lr{Lightning AI}.}
			\item {طراحی \lr{endpoint}های مقیاس‌پذیر برای دریافت \lr{prompt}، پردازش مدل و بازگردانی تصویر در سناریوهای تولید محتوای خودکار.}
			\item {یکپارچه‌سازی لایه سرویس با زیرساخت پردازش ابری به‌منظور تسهیل فرایند استقرار و بهره‌برداری از مدل.}
		\end{cvitems}
	}
	
	\cventry
	{اتوماسیون استخراج و ساخت جدول کلمات کلیدی از \lr{VidIQ}}
	{\lr{Python, Selenium, Web Automation}}
	{پروژه شخصی}
	{سال 1404}
	{
		\begin{cvitems}
			\item {توسعه اسکریپت اتوماسیون با \lr{Selenium} جهت استخراج کلمات کلیدی و متادیتاهای مرتبط از وب‌سایت \lr{VidIQ}.}
			\item {پالایش و تبدیل داده‌های استخراج‌شده به جداول ساختاریافته برای تحلیل \lr{keyword}ها و تدوین استراتژی محتوا.}
			\item {حذف کامل فرایندهای دستی جمع‌آوری داده و افزایش چشمگیر سرعت آماده‌سازی اطلاعات برای تحلیل محتوایی.}
		\end{cvitems}
	}
	
	\cventry
	{افزونه گزارش‌گیری زمانی \lr{(Time Tracking)} برای \lr{Google Calendar}}
	{\lr{Google Apps Script, Google Calendar, Google Sheets}}
	{پروژه شخصی}
	{سال 1403}
	{
		\begin{cvitems}
			\item {توسعه افزونه‌ای مبتنی بر \lr{Google Apps Script} برای پایش رویدادهای \lr{Google Calendar} و تولید خودکار گزارش در \lr{Google Sheets}.}
			\item {پیاده‌سازی سیستم گزارش‌گیری مبتنی بر \lr{tag} جهت دسته‌بندی فعالیت‌ها و تحلیل دقیق زمان تخصیص‌یافته به وظایف مختلف.}
			\item {خودکارسازی فرایند پایش زمان به‌منظور بهبود مستندسازی فعالیت‌ها و ارتقای بهره‌وری فردی.}
		\end{cvitems}
	}
	
	\cventry
	{اپلیکیشن اندروید دوست‌یابی بومی} % عنوان رسمی‌تر شد
	{\lr{Project Coordination, Mobile App Delivery, React Native}}
	{رایاکالج / \lr{Freelancer.com}}
	{سال 1401}
	{
		\begin{cvitems}
			\item {مدیریت و هماهنگی فازهای توسعه یک اپلیکیشن بومی در حوزه ارتباطات و تعامل کاربران (توسعه‌یافته با \lr{React Native}).}
			\item {نظارت بر روند اجرا، پایش نیازمندی‌ها و مدیریت پیشرفت پروژه تا مرحله استقرار و تحویل نهایی.}
			\item {تسهیل ارتباط با کارفرما و تضمین تطابق خروجی‌های توسعه با نیازمندی‌های کسب‌وکار.}
		\end{cvitems}
	}
	
	\cventry
	{اپلیکیشن تشخیص الگوهای صوتی برای افراد دارای محدودیت جسمی} % عنوان کمی شفاف‌تر شد
	{\lr{Java, Android Development}}
	{رایاکالج / \lr{Freelancer.com}}
	{سال ۱۳۹۹}
	{
		\begin{cvitems}
			\item {طراحی و برنامه‌نویسی اپلیکیشن اندروید جهت تشخیص الگوهای صوتی خاص و فعال‌سازی سیستم هشدار در شرایط اضطراری.}
			\item {پیاده‌سازی الگوریتم‌های پردازش ورودی صوتی با \lr{Java} و ادغام آن با مکانیزم‌های اطلاع‌رسانی سیستم‌عامل.}
			\item {ارائه راهکاری نرم‌افزاری با هدف تسهیل درخواست کمک برای افراد دارای محدودیت‌های حرکتی در موقعیت‌های بحرانی.}
		\end{cvitems}
	}
	
\end{cventries}


%---------------------------------------------------------
